\section[Diversité des publics]{Vers une bibliothèque numérique orientée usagers : s'adapter à la diversité des publics de Cujas}

La refonte de CujasNum ne saurait se limiter à une simple mise à niveau sur le plan technique. En plaçant l'usager au centre de son écosystème, cela représente un impératif stratégique important. La réussite de la future plateforme repose sur le passage d'un modèle sous forme de \enquote{répertoire de PDF} à un modèle de travail capable de répondre aux exigences techniques et méthodologiques permettant à un large public de l'utiliser, allant de l'étudiant de master au chercheur émérite. L'enjeu est de transformer l'infrastructure statique en un service dynamique soutenant la production scientifique.

L'analyse des entretiens et des retours d'enquête ont mis en évidence différentes typologies de recherche qui imposent des fonctionnalités différencies : l'enseignant-chercheur expert, le doctorant, l'étudiant et le lecteur éclairé. Ces quatre profils d'usager, utilisant la bibliothèque numérique, ont des méthodes de recherche propres. L'enseignant-chercheur est spécialiste d'une question ou d'une thématique, pratique une recherche de \enquote{fond} sur des sources rares. Ses habitudes de recherche incluent la sérendipité, l'utilisation de l'OCR comme d'un outil permettant de consulter le texte d'une façon plus agréable, et télécharge massivement des PDF pour alimenter ses prores dossiers. Le doctorant a des pratiques similaires à celui de l'enseignant-chercheur, mais exige davantage de précision sur la bibliothèque numérique utilisée. Il utilise le numérique pour la préparation pédagogique, nécessitnat des outils de capture d'écrans performant pour illustrer ses supports de cours. Pour l'étudiant, la bibliothèque numérique doit être efficace, passant par l'export bibliographique automatisé (Zotero), la rapidité d'accès, des explications précises et des guides d'aide à la prise en main de l'interface et de son contenu. Pour le lecteur éclairé, celui-ci va être attiré par les expositions virtuelles, ou une autre forme de valorisation sur des thématiques d'actualité sur lesquelles ils tomberaient par hasard pour en apprendre davantage.

Le numérique répond aujourd'hui à une tripe nécessité: logistique, scientifique et patrimoniale. Pour les chercheurs, la nouvelle bibliothèque numérique doit leur permettre d'avoir accès au contenu de la bibliothèque mais également provenant d'autres institutions françaises ou étrangères. Pour les institutions patrimoniales, l'usage du numérique représente un important levier de conservation. Face à la dégradation physique des collections papiées par le temps ou les intempéries, la version numérique offre une pérennité et une \enquote{propreté} que le document physique ne garantit plus. 
 
Actuellement, la bibliothèque numérique ne permet pas une adaptation optimale aux pratiques de la recherche, laissant présumé d'une compréhension par son public alors que ce n'est pas toujours le cas. Comme l'a expliqué Jérôme Dinet, en déconstruisant le mythe de l'autonomie des étudiants. Face aux développement des nouvelles technologies, il démontre que \enquote{l'aisance technique des nouvelles générations face aux écrans ne garantit pas leur capacité à chercher, filtrer et traiter de l'information académique}. Cette fracture d'usage entraîne des experts et des publics non-spécialiste légitime l'approche d'adapter la bibliothèque numérique en concevant des parcours différenciés et adaptés à chaque niveau de compétence. 
 
Par exemple, actuellement sur CujasNum, les retours des usagers face aux contenus proposés est un mélange d'incompréhension concernant la constitution des collections. En effet, sur le portail de CujasNum, on observe un \enquote{conflit de classification} entre des logiques thématiques et chronologiques. Il s'agit d'un flagrant contraste avec la promesse faite aux lecteurs lorsqu'il se trouve sur la page d'accueil, notamment avec les titres précis et les vignettes laissent entendre d'une forme de précision. Cette précision laisse entendre pour le lecteur qu'il trouvera un contenu spécialisé en consultant l'une des collections. Par exemple, dans la collection intitulée \enquote{Droit romain}, sans médiation éditoriale entraîne une confusion, à l'instar de la présence non-justifié pour les usagers d'auteurs du XVIIIe siècle, traitant du droit romain. Le retour fait notamment par les doctorants expliquent comme créant une collection \enquote{Droit romain}, ce n'est pas à cela qui s'attend mais plutôt à l'ensemble des lois, des règles et des principes juridiques régissant la Rome antique. Le manque d'explication derrière les choix faits par les experts métiers entraînent une confusion d'interprétation sur la logique suivie, engendrant une \enquote{peur de rater des sources}. Loin d'être confortable, c'est une des attentes majeures pour la nouvelle bibliothèque numérique de la part des usagers. Une autre de leurs attentes est de pouvoir retrouver le confort du livre papier qui est recréer numériquement. 